\section*{Appendix: Companion interfaces}
\addcontentsline{toc}{section}{Appendix: Companion interfaces}
\hypertarget{scale-law-interface}{}

This appendix records the scope of the estimates that can be used in
the companion limiting-law article \cite{CompanionLaw2026}. References below are to the numbered
statements of this paper. The finite-tree calculus, the asymptotic
comparison results, and uses of their proofs have different hypotheses;
these hypotheses remain in force under the normalizations described here.
Every number of labels, polynomial degree, derivative order, moment order,
and recursion depth is fixed before the system size tends to infinity.
Constants and the threshold in the size may depend on these fixed choices.

\paragraph{Fluctuation normalization.}
In \cite[Section~1]{CompanionLaw2026}, $s_n(\beta)^2$ denotes the
variance of the counting-prior, off-diagonal log partition function.
That free energy equals $F_n+n\log2$ in the notation of this paper,
so Theorem~\ref{thm:main} gives exactly
$s_n(\beta)=n^{1/6+o(1)}$. For the completed variable used in
hierarchical comparisons,
$\operatorname{Var}F_n^c=s_n(\beta)^2+T/2$.
The exponent theorem supplies neither a limiting law nor a variance
prefactor; those are separate conclusions of the companion.

\paragraph{Normalization and physical size.}
The hierarchy in Section~\ref{sec:foundations} has physical size $n$, spin
vector $y=\sigma/\sqrt n$, and terminal mass one. A hierarchy with mass
interval $[0,\ell]$, $1\le\ell\le4$, and terminal continuation
$\ell^{-1}\log(2^{-n}\sum_\sigma e^{\ell H(\sigma)})$ is obtained by
putting
\[
 u=v/\ell,\qquad
 \widetilde b(u)=\ell^2 b(\ell u),\qquad
 \widetilde h(u)=\ell^2 h(\ell u),\qquad
 \widetilde F=\ell F.
\]
The positive normalized transition laws are the same after this change
of variables. With self correction
$-\ell[b(\ell-)/4+h(\ell-)/2]$, its pressure satisfies
$f_n^{[0,\ell]}(b,h)=\ell^{-1}f_n(\widetilde b,\widetilde h)$.
Conditional means, covariances and overlaps are unchanged; the
normalized-source Hessians satisfy $H_v=\ell\widetilde H_{v/\ell}$.
Consequently $vC_v\preceq H_v\preceq\ell C_v$, and the pressure
formula retains the coefficients $-1/4$ and $-1/2$ when integrated in
$dv$. Terminal coincidence weights in the signed calculus become
$\ell$, while edge and fork weights are $-dv$ and $-(d-1)v$.

An ambient dyadic parameter $N$ may be used when $N/2\le n\le3N$;
this changes fixed power or logarithmic scale comparisons only by
bounded factors. The physical $n$ remains in Gaussian covariances,
overlap denominators and innovation estimates. On deleting $j$ fixed
sites, keep the original denominator $n$ in the deleted overlap and
keep all undeleted Gaussian covariances unchanged. If the retained
system is instead written with its own size $p=n-j$ and overlap
$R_p$, its matrix clock is $(p/n)b$ and its field clock is $h$.
Replacing a deleted overlap by the full one costs at most $j/n$;
derivatives and finite differences retain the associated weights and
self corrections. A size stencil at fixed clocks uses the same named
paths, cuts and external scale parameters in every stencil term.
Comparability of $n$ and $N$ does not relax a direct-test cutoff or
remove the fixed positive exponent slack in \eqref{eq:3}.

\paragraph{Source and pressure calculus; the scope of pinned moments.}
\Cref{lem:found-source,lem:found-pressure,cor:found-source-bounds}
give the source derivatives, pressure derivatives and fixed-order
source bounds for the normalized positive hierarchies. A common
deterministic source is permitted in the source and pressure
statements. The spin moment inequality of
\Cref{lem:found-moments}, however, is used at zero deterministic
external field, with deterministic weights, under the full
unconditioned single-path planted law. Its centering is the expectation
under that law. It is not a pinned inequality for an arbitrarily
conditioned prefix or a spin-dependent reweighting of that marginal.

The same proof gives the following useful polynomial form. Let $P_u$
have a fixed degree and bounded deterministic coefficients, measurable
in $u$, and let $k\in L^2$ be deterministic. At zero deterministic field,
\[
 \left\|\int k_u\bigl(P_u(L_u)-\mathbf E P_u(L_u)\bigr)\,du\right\|_p^2
 \le c_p\left\|\int k_u^2P'_u(L_u)^2W_u\,du\right\|_{p/2},
 \qquad 2\le p<\infty.
\]
Here $c_2=1$ and one may use the same constants $c_p\le1+p^2/3$.
Indeed, in the proof of \Cref{lem:found-moments}, the driver gradient
of $P_u(L_u)$ is $P'_u(L_u)g_u$. The positive block-Hessian argument
therefore gives the same inverse-Hessian energy bound, with
$J=\int k_u^2P'_u(L_u)^2W_u\,du$. Coordinate gauge symmetry preserves
both the observable and this energy and makes their conditional laws
independent of the pinned spin. The same Brascamp--Lieb and
real-moment argument applies. Boundedness of $L_u$ passes the result
from bounded deterministic weights to $L^2$ weights. The preceding
mass rescaling gives the identical formula on $[0,\ell]$.

For a fixed positive enlarged tree, first marginalize all branches
unused by a given single-path factor. Apply the single-path inequality
to that full marginal; next combine the retained factors by
unconditional H\"older inequalities under the positive tree law; and
only then integrate against the total variation of the deterministic
signed allocation measure. This is the order in
\Cref{rem:found-single-pin,comp:good-split-trees}. Random factors on
other branches are not substituted for deterministic weights in the
single-path inequality. No joint-spin pinned log-concavity assertion
is required by this procedure.

\paragraph{Projection, resampling and the conditioning field.}
\Cref{comp:projection} applies when, conditional on the data determining
the exterior anchor and the shared prefix, the selected branch retains
its entire normalized planted future law. The anchor may use an
independent outside continuation after the fork. It may not use an
unrevealed variable from the selected continuation. After a field-clock
interval of length $\Delta h$, at masses at least $v>0$, the mass-$\ell$
normalization gives
\[
 \|V\cdot(y-X_{\rm end})\|_p
 \le C_p\frac{(\log n)^2\|V\|_p}{v\sqrt{n\Delta h}},
\]
for fixed even $p$, with mass and length bounded below by fixed inverse
powers of $n$. Conditional Jensen permits later projections; independent
conditional resampling uses two such replacements. The proof also applies when total clocks grow by a fixed power of
$\log n$, as in the companion application. More generally, polynomially
bounded field intervals leave the initial-to-innovation ratio in its
moment iteration polynomially bounded. The normalized innovation
estimate is unchanged. This extension concerns the projection estimate
itself, not the asymptotic comparison theorem.

For \Cref{lem:found-coupling}, the conditioning field $\mathcal G$ is
\emph{exactly} the sigma-field generated by the complete common prefix
through the specified revelation at the cut, the old-branch
realizations outside the affected descendant subtree after that cut,
and the deterministic old topology. It includes no information from
the new row's continuation or the affected old subtree after the cut.
The two conditional marginals are then the same under the two laws.
The centered-product estimate uses unconditional $L^p$ norms after
averaging over this $\mathcal G$; no bound uniform in the realized
prefix is asserted. Preserving a conditional mean alone would not
preserve the hypotheses of the projection or pinned-moment estimates.

\paragraph{Signed allocations and fixed-size path limits.}
\Cref{prop:found-replica,lem:found-attachment,lem:found-split-density,lem:found-taylor}
and \eqref{eq:2} concern ordinary positive tree laws integrated against
signed deterministic coefficients. All fork atoms and terminal
coincidences are retained. The bounds have no mesh or minimum-mass
loss for a fixed number of labels. Allocation-order symmetry carries
the complete topology coefficient. A row-sum cancellation is made only
when the remaining integrand is independent of that row's detailed
placement. In a recursive calculation, freeze the previously generated
topology, its coefficients and its scalar centers before introducing
the next fresh row. The restrictions in
\Cref{lem:found-cut-truncation,lem:found-boundary-tests} remain part of
the interface.

The Taylor and bounded-row-change estimates require a uniformly
bounded covariance difference along a positive semidefinite
interpolation, even when an individual clock is long. In a comparison
of one matrix row with a scalar row at $K$, it is the row matrix
coefficient and the difference $h-K$ that must obey the stated bounds.
Clock-independent identities and fixed-size limits alone do not
give a bounded-factor row transfer for an unbounded covariance change.
If a test, its scalar centers or its deterministic weights vary with
the interpolation, their product-rule derivatives must also be kept.

\Cref{prop:found-limits} is a fixed-$n$ limit statement. Terminal clocks
and the specified before/after clock values at each retained cut must
converge. At a jump of the pair $(b,h)$ use the fixed joint filling
\[
 (b_-+t\Delta b,h_-+t\Delta h),\qquad 0\le t\le1,
\]
at the jump's fixed mass. Intermediate cuts retain their paired clock
values in this filling. Positive tree expectations at unfilled cuts,
fixed-order source derivatives, signed allocations, and the
mass-integrated pressure and moment formulas pass to the limit as
specified in that proposition. The susceptibility inequality holds
on continuous parts, with a nonnegative total increment of $S$ at
each jump. To retain separate driver contributions or intermediate
cuts inside a jump, the approximations must retain the chosen joint
filling. The total increment of $S$ does not depend on this filling,
but its two driver allocations need not agree for different fillings.
No discretization rate uniform in the physical size is inferred.

\paragraph{Finite-degree cavity calculations under changed constraints.}
The named proof modules in Section~\ref{sec:cavity} use the quantitative
analytic inputs listed there: raw moments, normalized-future
projection and resampling, a deterministic good split with a saving
independent of each fixed moment order, integrated high-split
estimates with the required split densities, and the low-scale
integrated and target errors. The scalar comparison uses the entire
path $K=h+bS$, not only its value at the target. Bounded one-row changes
preserve positive semidefiniteness and leave the other row covariances
unchanged; the row transfer and deletion errors in
\Cref{lem:cavity-row-change} are paid at their actual scales.

Under these analytic inputs and the test-family restrictions of
\Cref{def:cavity-tests,lem:cavity-family}, the row identity, projected
replacements and fixed-depth cancellation in
\Cref{lem:cavity-row-identity,lem:cavity-replacements,lem:cavity-recursion}
are available with their displayed remainders. In each generation the
current scalar coefficients and centers stay fixed during its new
interpolation. Observed-block rows, fork atoms, and factors with a
coefficient depending on the current placement must be processed as
in those proofs, before any allowed row collapse or absolute bound.
Each terminal marked star factor contributes its saved fixed moment;
one chooses a finite depth sufficient for the desired power margin,
then fixes all resulting degrees and H\"older orders before taking
$n\to\infty$.

Changing the inverse-density constraints does not by itself give the
conclusion of \Cref{prop:cavity-closure,prof:direct-profile} for the
new optimizer. Their proofs can be used after establishing, for that
optimizer, attainment and admissibility, the appropriate first
variation and feasible competitors, the analytic inputs just listed,
and the required mass-\emph{and}-radius error estimates. In particular,
small floor mass does not imply small floor radius, and a whole-cap
trace estimate is not a pointwise cap estimate. The endpoint expansion
used for susceptibility calibration is required on the regular
calibration bands; no such expansion is silently imposed at an
omitted lower target. These are proof obligations for a modified
comparison, not consequences of having a similar objective.
Likewise, a general finite cavity quotient on other physical paths
must be derived from the positive covariance interpolation and finite
Taylor calculus at those paths; it is not an unrestricted instance of
the minimizing-path closure proposition.

\paragraph{The finite-propagation output.}
\Cref{reg:finite-propagation} is used with the actual family of
\Cref{reg:families}, the taper in \eqref{eq:6}, and the constraints
in \eqref{eq:5}. In particular, for a sufficiently small \emph{fixed}
$m>0$ and $\delta=m/10000$,
\[
 n^{-2/3+m}d\lesssim E\le n^{-1/2}d^2,\qquad
 a=\sqrt n E/d,\qquad u_*=(nE)^{-1},\qquad 0<d=O(1),
\]
and the sequential endpoint condition is
\[
 b_0(u)=T+o(z^2),\qquad h_0(u)=o(z^3)
 \quad (u\asymp z,\ a\lesssim z\le1).
\]
Total clocks and all fixed comparison constants are uniformly bounded;
$d^2\le T-\rho$ for a fixed $\rho>0$. All changed feasible paths
are nonnegative and nondecreasing. The global matrix lower bound or
the high-only matrix lower bound and boundary gap has a fixed positive
margin after the initial change, available for the entire finite
descendant chain. An omitted mass interval satisfies $u_0=o(a)$,
and its wall satisfies the stated constraint in \eqref{eq:5}.

Under these assumptions every minimizing quantile obeys
\[
 cu\le r(u)\le Cu
 \quad\left(n^\delta\sqrt n E/\sqrt\lambda\le u\le1\right),
 \qquad 0<\lambda\le d^2.
\]
For an ordinary family with an admissible bare endpoint, its minimized
excess lies between zero and $n^{20\delta}E$ throughout the sweep.
For a fixed-strength family only the profile conclusion is asserted;
an inadmissible bare root is not evaluated. Before finite propagation
has supplied regularity, \Cref{comp:conditional-error} requires an
outer regular region above $A_\lambda$ when $A_\lambda=o(1)$.
Likewise a small direct test requires the outer regularity and all
three local inequalities stated in \Cref{prof:direct-profile}.
The propagation proof obtains the child's outer region from
successful-error agreement before applying that direct result.

\paragraph{A common prelimit temperature.}
The same finite-propagation and direct-test conclusions hold
sequentially if the endpoint condition uses $T_n\to T_\infty>1$,
with $T_n$ held fixed throughout the whole parent--child construction
at size $n$, and with all the preceding fixed margins and bounds.
Here is the required proof-level replacement. At the regular
calibration cut retain
\[
 b_*=T_n+O(W^2),\qquad K_*=T_nS_*+O(W^3),\qquad
 J_0=T_n\chi^2-1
\]
in \eqref{eq:26}--\eqref{eq:27}. Division uses
$\chi K_*\asymp W$ and the same absorption of the
$\eta_n|J_0|/W^2$ term. No error $T_n-T_\infty$ is introduced.
In the scalar compactness and linearization arguments use time
$K/T_n$, subtract $K/T_n$, and keep
$[(b_0-T_n)S+h_0]/x^3$ as the endpoint remainder. The derivative,
strict-curvature and compactness bounds use a fixed compact
neighborhood of $T_\infty$ and the total-clock bounds. Only in the
limiting equations is $T_n$ replaced by $T_\infty$.

At a macroscopic direct test the endpoint convergence gives the
limiting variational problem at $T_\infty$; the scalar-minimizer
result \Cref{prof:scalar-minimizer} is used at that one fixed
temperature. The child-data estimates retain the same $T_n$ in
each endpoint expansion, and covariance margins are chosen with room
at the limiting temperature. The same fixed-depth contradiction
therefore applies. This sequential extension requires no rate for
$T_n-T_\infty$, gives no uniform assertion as the limiting temperature
approaches one, and permits neither growing degree nor an
$n$-dependent recursion depth.
